The descent into madness from sleep deprivation provides a terrifying insight into sleep's function, especially when combined with stimulants. Amphetamine-induced psychosis, common in stimulant addiction, is clinically similar to schizophrenia and is exacerbated by the extended wakefulness the drugs enable \citep{lechi2012amphetamine}.

This suggests a "full-system overfitting failure," a "Tetris effect gone rogue."
In the Tetris effect, the brain "overfits" to a neutral visual pattern (falling blocks). In stimulant- or deprivation-induced psychosis, the brain—flooded with dopamine (the "salience" chemical) and deprived of its nightly "reset"—overfits to \textit{internal} patterns of fear, vigilance, and threat.

Without the NREM phase to "downscale" and prune these circuits, and without the REM phase to "regularize" and "diversify" the brain's model of the world, a paranoid feedback loop begins. The brain gets "stuck," seeing threat patterns everywhere, and the mind loses its grip on a shared reality. The fact that these severe psychotic symptoms often resolve within a day or two of "crashing" and getting recovery sleep further reinforces that sleep is the brain's essential, non-negotiable error-correction and maintenance phase \citep{waters2018deprivation}.